\documentclass[11pt]{article}
\usepackage[margin=0.75in]{geometry}
\usepackage{amsmath,amssymb,booktabs,graphicx,hyperref}
\newcommand{\todo}[1]{\textbf{[TODO: #1]}}

\title{T1: Stallion Biomechanics, Checkpoint 2: Implementation and Evaluation}
\author{Jim Daou, Stephen Gancayco, Kelly Tran, Samuel Whitney\\
CPSC 4430/6430, Fall 2026}
\date{}

\begin{document}
\maketitle
\vspace{-1.5em}

\section{Implementation}

We implemented the Checkpoint 1 plan end to end. PosePAL~\cite{posepal}
labels points in a video without training a species-specific detector:
a general-purpose point tracker, CoTracker3~\cite{cotracker3}, stays
frozen, and only a per-keypoint appearance embedding (one vector per
landmark, \texttt{kp\_feats} in the authors' code) is optimized for each
video at test time. The embedding is initialized from frame 1 and the
annotated frames, then updated with Adam on the tracker's own loss,
$L = L_{\text{track}} + 0.01 \cdot L_{\text{reg}}$, where
$L_{\text{track}}$ is the CoTracker3 tracking loss (Huber, $\delta = 6$,
$\gamma = 0.8$, over refinement iterations) and $L_{\text{reg}}$ is an
L1 penalty toward the initialization.

Our pipeline has four stages. (1) Sparse manual annotation of the 28
custom landmarks (24 horse, 4 dummy) as described in Section 2.
(2) Tracking: a frozen CoTracker3 baseline, and PosePAL test-time
optimization (TTO) using the correction frames as supervision.
\todo{Stephen: state the TTO settings actually used, steps and learning
rate schedule; the paper uses 1{,}000 steps at $10^{-3}$ to $10^{-5}$,
the released code defaults to 500 steps at $10^{-4}$}
(3) Endpoint computation from the dense tracks (Section 2).
(4) Phase prediction over approach, mount, thrust, and dismount, with a
rule-based predictor over the endpoint series and a random forest
(pose-only features vs.\ temporal features). Group code (annotation
tooling, endpoint and phase modules, evaluation, demo renderer) is in
the project repository, Section 6; the tracker itself is the authors'
released code at commit 63572cd.

\section{Data and splits}

\textbf{Stallion clips.} \todo{Sam: clip table, clip IDs, horse IDs,
lengths, source; the demo clip must be named} The Checkpoint 1
commitment is at least 100 reference frames across at least 5 clips;
we achieved \todo{Sam: actual counts}.

\textbf{Paper datasets.} The tracking reproduction also runs on
DAVIS-Animals~\cite{davis} and DeepFly3D~\cite{deepfly3d}, the paper's
evaluation data, so our frozen-vs-adapted comparison can be sanity
checked against published numbers (paper $\delta_{\text{avg}}$ 67.53 on
DAVIS-Animals; those numbers do not transfer to our 28 custom
landmarks, for which no pseudo-labels exist and all annotation is
manual).

\textbf{Annotation protocol.} For each clip, all 28 landmarks are
hand-marked on frame 1 (the tracker query points) and on every
$\sim$10th frame (correction frames, fed to TTO). The interleaved
frames (5, 15, 25, \dots) are annotated in a separate file as blind
scoring frames and are never used for correction or model selection.
Splits are by horse when possible, otherwise by whole clip, keeping
temporal windows together, so no frame of a test clip influences its
own adaptation. Occluded landmarks are marked invisible and reported
missing; nothing is interpolated in annotation, endpoints, or scoring.

\textbf{Endpoints (features).} Four image-plane quantities per frame,
with Checkpoint 1 sign conventions: back angle
$\theta_{\text{back}}$, the withers-to-croup line against horizontal,
positive when the croup is above the withers; neck angle
$\theta_{\text{neck}}$, the poll-to-withers line against the back-line
direction, positive when the poll is above the back line; hind-limb
position, the horizontal hoof offset from the pelvis vertical divided
by torso length $L_{\text{torso}}$ (withers to croup), positive
forward; and chest position, the horizontal chest offset from dummy
landmark D1 divided by $L_{\text{dummy}}$ (D1 to D2), positive toward
the dummy. An endpoint is computed only when all of its landmarks are
visible that frame.

\section{Evaluation protocol}

Metrics follow the paper and the Checkpoint 1 precommitment. Tracking:
$\delta_{\text{avg}}$, the fraction of points within 1, 2, 4, 8, and 16
px at $256 \times 256$; $J$, the Jaccard-style tracking quality; and
$J_{\text{masked}}$, which weights large errors (at least 4 px) ten
times more. \textbf{The declared primary metric is $J_{\text{masked}}$
on the blind held-out frames.} Endpoints: mean absolute error per
endpoint against the manual reference on scoring frames, plus coverage
(the fraction of frames where the endpoint is computable). Phases:
macro-F1 over the four phases. Model choices (tracker settings,
endpoint rules, phase model class) were frozen before final testing:
\todo{group: date frozen}. Declared success target, fixed before final
testing: TTO improves $J_{\text{masked}}$ over the frozen baseline on
held-out frames \todo{group: add the numeric targets for endpoints and
phase macro-F1 before running final tests}. Improvement experiment
(Checkpoint 1 hypothesis, fixed annotation budget): correction frames
selected by tracker uncertainty (low visibility confidence, expected
to concentrate in mount and thrust) vs.\ the uniform every-10th-frame
schedule, with the same number of annotated frames; only the positions
change.

\section{Results}

\todo{Jim: fill Tables 1 to 3 from the frozen evaluation run; one
sentence of interpretation per table; add the demo overlay still as
Figure 1 if space allows}

\begin{table}[h]
\centering\small
\begin{tabular}{lccc}
\toprule
 & $\delta_{\text{avg}}$ & $J$ & $J_{\text{masked}}$ (primary)\\
\midrule
Frozen CoTracker3 & TBD & TBD & TBD\\
PosePAL TTO (uniform corrections) & TBD & TBD & TBD\\
PosePAL TTO (uncertainty-selected) & TBD & TBD & TBD\\
\bottomrule
\end{tabular}
\caption{Tracking on blind held-out frames. \todo{Jim: state the clip
set and frame counts}}
\end{table}

\begin{table}[h]
\centering\small
\begin{tabular}{lcc}
\toprule
Endpoint & MAE & Coverage\\
\midrule
$\theta_{\text{back}}$ (deg) & TBD & TBD\\
$\theta_{\text{neck}}$ (deg) & TBD & TBD\\
Hind-limb position (norm.) & TBD & TBD\\
Chest position (norm.) & TBD & TBD\\
\bottomrule
\end{tabular}
\caption{Endpoint errors against the manual reference. \todo{Sam/Jim}}
\end{table}

\begin{table}[h]
\centering\small
\begin{tabular}{lc}
\toprule
Phase model & Macro-F1\\
\midrule
Rule-based (endpoint thresholds) & TBD\\
Random forest, pose-only & TBD\\
Random forest, temporal features & TBD\\
\bottomrule
\end{tabular}
\caption{Phase prediction over approach, mount, thrust, dismount.
\todo{Jim: add per-phase F1 if space allows}}
\end{table}

\section{Limitations}

Endpoints are image-plane quantities; without camera calibration they
are not physical angles or distances. Coverage is lowest exactly where
the biology is most interesting: in mount and thrust, landmarks are
occluded by the dummy and the stallion's own body, and our missing-data
rule (report, never interpolate) turns that occlusion into missing
endpoints rather than noisy ones. All supervision is manual, so the
annotation budget caps both clip count and correction density.
\todo{group: honest target assessment, which declared targets were met
or missed and what was run in response; the assignment requires this}

\section{Reproducibility}

Repository: \url{https://github.com/kellytran969/T1-Stallion-Biomechanics},
submitted revision \todo{commit hash or tag at submission}. The README
records setup, data and model access, exact commands, expected outputs,
and limitations. Authors' code: PosePAL repository at commit 63572cd
with the CoTracker3 submodule at 493d363; tracker weights
\texttt{scaled\_offline.pth} (97 MB) from Hugging Face. Hardware:
\todo{Stephen: machine and GPU used, runtime per clip; paper reference
is about 3 minutes per 100-frame 480p clip on an RTX A6000}.

\begin{thebibliography}{9}\small
\bibitem{posepal} Z. Pan, B. Pan, G. Yang, A. W. Harley, L. Guibas.
Animal Pose Labeling Using General-Purpose Point Trackers.
arXiv:2506.03868, 2025.
\bibitem{cotracker3} N. Karaev, Y. Makarov, J. Wang, N. Neverova,
A. Vedaldi, C. Rupprecht. CoTracker3: Simpler and Better Point Tracking
by Pseudo-Labelling Real Videos. arXiv preprint, 2024.
\bibitem{davis} J. Pont-Tuset, F. Perazzi, S. Caelles, P. Arbel\'aez,
A. Sorkine-Hornung, L. Van Gool. The 2017 DAVIS Challenge on Video
Object Segmentation. arXiv preprint, 2017.
\bibitem{deepfly3d} S. G\"unel, H. Rhodin, D. Morales, J. Campagnolo,
P. Ramdya, P. Fua. DeepFly3D, a deep learning-based approach for 3D
limb and appendage tracking in tethered, adult Drosophila.
eLife, 2019.
\end{thebibliography}

\section*{Contributions}
\small
\begin{tabular}{ll}
\toprule
Member & Checkpoint 2 parts\\
\midrule
Stephen Gancayco & Implementation lead: PosePAL runs (frozen and TTO), repository\\
Samuel Whitney & Data, split, endpoint measurements\\
Jim Daou & Evaluation, comparisons, phase prediction\\
Kelly Tran & Annotations, demo video, README and report package\\
\bottomrule
\end{tabular}

\end{document}
